% MemQ v14 -- joint WebQSP + CWQ + GrailQA fine-tuning: result tables.
%
% Standalone-compilable (pdflatex report.tex); the individual table
% environments can also be copied into an existing report.
% All figures are measured, not estimated; the provenance of each number is
% given in the accompanying notes.
\documentclass[11pt,a4paper]{article}
\usepackage[margin=2.5cm]{geometry}
\usepackage{booktabs}
\usepackage{siunitx}
\usepackage[hidelinks]{hyperref}
\sisetup{detect-weight=true,detect-family=true}
\sloppy

\title{MemQ v14: Joint Fine-tuning on WebQSP, CWQ and GrailQA}
\author{}
\date{}

\begin{document}
\maketitle

\section{Data composition}

The v14 model is fine-tuned jointly on the training splits of all three
benchmarks. GrailQA dominates the corpus because its training split is an order
of magnitude larger than WebQSP's. GrailQA's official test split is unlabelled
(leaderboard only), so its \emph{dev} split serves as the evaluation set, as is
common practice.

\begin{table}[htbp]
  \centering
  \caption{Data used for training and evaluation. No test item appears in the
  training corpus (verified by identifier intersection).}
  \label{tab:composition}
  \small
  \begin{tabular}{lS[table-format=5.0]S[table-format=3.1]lS[table-format=5.0]}
    \toprule
    & \multicolumn{2}{c}{Training} & \multicolumn{2}{c}{Evaluation} \\
    \cmidrule(lr){2-3}\cmidrule(lr){4-5}
    Dataset & {Items} & {Share (\%)} & Split & {Items} \\
    \midrule
    GrailQA & 44337 & 59.2 & dev  & 6763 \\
    CWQ     & 27474 & 36.7 & test & 3531 \\
    WebQSP  &  3085 &  4.1 & test & 1637 \\
    \midrule
    Total   & 74896 & 100.0 & & 11931 \\
    \bottomrule
  \end{tabular}
\end{table}

Table~\ref{tab:share} contrasts the corpus with that of the earlier v9 model.
Although five epochs expose every WebQSP example exactly as often as in v9, its
share of any individual gradient step drops by more than half, while 44\,337
GrailQA examples compete for capacity.

\begin{table}[htbp]
  \centering
  \caption{Share of the training corpus per benchmark. v9 was trained on
  WebQSP and CWQ only.}
  \label{tab:share}
  \begin{tabular}{lS[table-format=2.1]S[table-format=2.1]}
    \toprule
    Dataset & {v9 (\%)} & {v14 (\%)} \\
    \midrule
    WebQSP  & 10.1 &  4.1 \\
    CWQ     & 89.9 & 36.7 \\
    GrailQA & {--} & 59.2 \\
    \bottomrule
  \end{tabular}
\end{table}

The leakage check is not vacuous: CWQ derives its questions from WebQuestions
and reuses the base identifier with an appended hash
(\texttt{WebQTest-322\_bd42258\ldots}), so \num{8493} training items carry a
\texttt{WebQTest} prefix without being WebQSP test items. Intersection with
WebQSP test, CWQ test and GrailQA dev is empty in all three cases.

\section{Training configuration}

The official MemQ repository prescribes only ``follow the LLaMA-Factory,
\texttt{lora rank=16}, \texttt{epoch=5.0}''; the paper reports no fine-tuning
hyperparameters at all. The v9 configuration is therefore the reference for
comparability. Sequence packing is the one deliberate deviation: without it,
per-batch padding produces varying GEMM shapes that intermittently trigger an
illegal memory access in a Tensile kernel on gfx1201, which terminates the run.
Packing fixes every batch at \(2 \times 2048\) tokens and additionally removes a
large amount of padding waste.

\begin{table}[htbp]
  \centering
  \caption{Training configuration. The number of optimiser updates is kept close
  to v9 so that the learning rate carries over unchanged.}
  \label{tab:config}
  \begin{tabular}{lll}
    \toprule
    Parameter & v9 & v14 \\
    \midrule
    Base model            & \multicolumn{2}{l}{Llama-3-8B-Instruct (bfloat16)} \\
    LoRA rank / target    & \multicolumn{2}{l}{16 / all linear} \\
    Learning rate         & \multicolumn{2}{l}{\num{5e-5}, cosine, warmup 0.1} \\
    Epochs                & \multicolumn{2}{l}{5.0} \\
    Cutoff length         & \multicolumn{2}{l}{2048} \\
    \midrule
    Packing               & no            & yes (neat) \\
    Batch $\times$ accum. & $2 \times 4$  & $2 \times 1$ \\
    Optimiser updates     & 18145         & 17390 \\
    Trainable parameters  & \multicolumn{2}{l}{\num{41943040} (0.52\%)} \\
    Wall-clock time       & {--}          & \SI{20.8}{\hour} \\
    \bottomrule
  \end{tabular}
\end{table}

\section{Checkpoint selection}

The run recorded no validation loss (\texttt{eval\_strategy: "no"}), so the five
per cent held-out split was scored retrospectively for seven candidates
(Table~\ref{tab:checkpoints}). The earliest retained checkpoint has the lowest
loss and the final one is \SI{5.8}{\percent} worse, which suggests mild
overfitting.

That ranking does not survive contact with the benchmarks. Table~\ref{tab:cand}
shows the opposite order: the checkpoint with the \emph{worst} held-out loss
produces the \emph{best} WebQSP answers, and the loss-optimal one is last on
both benchmarks. Held-out loss is therefore not a usable selection criterion for
this task; cp15000 was chosen because it leads on GrailQA, the target benchmark.

\begin{table}[htbp]
  \centering
  \caption{Held-out loss per checkpoint, on the identical five per cent split
  (\num{348} packed sequences).}
  \label{tab:checkpoints}
  \begin{tabular}{lS[table-format=5.0]S[table-format=1.4]S[table-format=+1.1]}
    \toprule
    Checkpoint & {Step} & {Held-out loss} & {vs.\ best (\%)} \\
    \midrule
    cp13750 & 13750 & 0.0104 & 0.0 \\
    cp14250 & 14250 & 0.0105 & +1.0 \\
    cp14000 & 14000 & 0.0107 & +2.9 \\
    cp15000 & 15000 & 0.0108 & +3.8 \\
    cp14500 & 14500 & 0.0109 & +4.8 \\
    cp14750 & 14750 & 0.0110 & +5.8 \\
    final   & 17390 & 0.0110 & +5.8 \\
    \bottomrule
  \end{tabular}
\end{table}

\begin{table}[htbp]
  \centering
  \caption{Candidate comparison. WebQSP: \num{300} test questions; GrailQA:
  the \num{200}-question sample from the v9 transfer experiment. Paired
  comparison, identical questions for all candidates.}
  \label{tab:cand}
  \begin{tabular}{lS[table-format=1.4]S[table-format=1.4]S[table-format=1.4]}
    \toprule
    Checkpoint & {Held-out loss} & {WebQSP F1} & {GrailQA F1} \\
    \midrule
    cp13750 & \bfseries 0.0104 & 0.7991 & 0.2704 \\
    cp15000 & 0.0108 & 0.8114 & \bfseries 0.2826 \\
    final   & 0.0110 & \bfseries 0.8129 & 0.2751 \\
    \bottomrule
  \end{tabular}
\end{table}

\section{Transfer to GrailQA}

This is the central result. v9, trained on WebQSP and CWQ only, does not
transfer: on the \num{200}-question sample it matches not a single gold relation
multiset and its edge hitting rate is below one per cent (Hits@1 \num{0.040},
F1 \num{0.037}). v14 reaches \num{0.285}/\num{0.283} on the same sample. The
full dev set (\num{6763} questions) is evaluated below; it also exercises the
operators, which the sample did not (it contained only \texttt{function:
none}).

\begin{table}[htbp]
  \centering
  \caption{GrailQA dev, \num{6763} questions, v14 (cp15000). The headline number
  is dominated by zero-shot questions, which are 54\,\% of the set.}
  \label{tab:grailqa}
  \begin{tabular}{lS[table-format=4.0]S[table-format=1.4]S[table-format=1.4]}
    \toprule
    Split & {n} & {Macro-F1} & {Hits@1} \\
    \midrule
    overall             & 6763 & 0.3355 & 0.3305 \\
    \midrule
    i.i.d.              & 1593 & \bfseries 0.8185 & 0.8129 \\
    compositional       & 1514 & 0.4266 & 0.4287 \\
    zero-shot           & 3656 & 0.0873 & 0.0796 \\
    \bottomrule
  \end{tabular}
\end{table}

On i.i.d.\ questions v14 matches its own WebQSP level (\(F_1 = 0.82\)); the
overall figure is pulled down by the zero-shot majority. The cause is
structural, not a defect: MemQ reconstructs a query by retrieving relations from
a memory built out of the \emph{training} gold queries. A zero-shot question by
definition needs unseen relations, and Table~\ref{tab:coverage} shows the
memory simply does not contain them -- so the correct relation cannot be
retrieved regardless of plan quality. This is an inherent limit of a
memory-based method, and it bounds what any amount of tuning can achieve on the
zero-shot split.

\begin{table}[htbp]
  \centering
  \caption{Fraction of each dev split's gold relations that exist in the v14
  memory. The memory is built from training queries, so it covers i.i.d.\ and
  compositional questions completely but almost none of the zero-shot ones.}
  \label{tab:coverage}
  \begin{tabular}{lS[table-format=3.1]}
    \toprule
    Split & {Gold relations in memory (\%)} \\
    \midrule
    i.i.d.        & 100.0 \\
    compositional & 100.0 \\
    zero-shot     &   4.5 \\
    \bottomrule
  \end{tabular}
\end{table}

\paragraph{The operators work.} Table~\ref{tab:ops} breaks the dev set down by
function. The \num{1251} operator questions reach \(F_1 = 0.41\), \emph{above}
the \num{0.32} of the plain questions, so the extension is not merely nominal.
The reconstructed queries carry the right construct:
\SI{96}{\percent} of argmax/argmin plans emit \texttt{ORDER BY} and
\SI{99}{\percent} of count plans emit \texttt{COUNT()}. This is the direct
contrast to CWQ superlatives (Section~\ref{sec:defects}, ``superlatives never
learned''), where the same construct is produced \SI{0}{\percent} of the time --
there the operator is absent from the training data, here it is present.

\begin{table}[htbp]
  \centering
  \caption{GrailQA dev by operator, v14 (cp15000).}
  \label{tab:ops}
  \begin{tabular}{lS[table-format=4.0]S[table-format=1.4]S[table-format=1.4]}
    \toprule
    Function & {n} & {Macro-F1} & {Hits@1} \\
    \midrule
    none (path only)     & 5512 & 0.3182 & 0.3139 \\
    \midrule
    count                &  440 & 0.3205 & 0.3205 \\
    argmax               &  298 & 0.4450 & 0.4698 \\
    argmin               &  210 & 0.4381 & 0.4381 \\
    comparison ($<\le>\ge$) & 303 & 0.4921 & 0.4356 \\
    \midrule
    all operators        & 1251 & \bfseries 0.4115 & 0.4037 \\
    \bottomrule
  \end{tabular}
\end{table}

\section{WebQSP: cost of joint training}

\begin{table}[htbp]
  \centering
  \caption{WebQSP, \num{300} questions, paired comparison. v14 builds
  structurally better queries but answers slightly less accurately.}
  \label{tab:webqsp}
  \begin{tabular}{lS[table-format=1.4]S[table-format=1.4]S[table-format=+2.1]}
    \toprule
    Metric & {v9} & {v14 (cp15000)} & {$\Delta$ (pp)} \\
    \midrule
    Structure accuracy   & 0.4500 & \bfseries 0.6367 & +18.7 \\
    Hits@1               & \bfseries 0.8423 & 0.8188 & -2.4 \\
    Macro-F1             & \bfseries 0.8542 & 0.8114 & -4.3 \\
    \bottomrule
  \end{tabular}
\end{table}

This reproduces the divergence already reported for v11, which improved
structure accuracy substantially while regressing answer-level metrics. Two
explanations are consistent with the data and cannot be separated here: the
DeepSeek-generated memory phrasing, and the reduced WebQSP share of the joint
corpus (Table~\ref{tab:share}). Note that v14 uses the same memory at training
and inference time, so train/inference consistency alone does not prevent the
effect.

\section{Reproduction of the paper and the remaining gap}

Table~\ref{tab:repro} compares the v9 baseline, evaluated over the full test
sets, against the paper. CWQ reproduces the previously documented figures
almost exactly (0.7241/0.7385 versus 0.725/0.738), which validates the
measurement chain; the gap to the paper is therefore a property of this
reproduction, not of the evaluation code.

\begin{table}[htbp]
  \centering
  \caption{v9 over the full test sets versus the paper (Table 3 of the
  original). EHR is reported after the fix described in
  Section~\ref{sec:defects}.}
  \label{tab:repro}
  \small
  \begin{tabular}{llS[table-format=1.4]S[table-format=1.3]S[table-format=+2.1]}
    \toprule
    Dataset & Metric & {This work} & {Paper} & {$\Delta$ (pp)} \\
    \midrule
    WebQSP & Hits@1   & 0.8229 & 0.857 & -3.4 \\
           & Macro-F1 & 0.8237 & 0.872 & -4.8 \\
           & EHR      & 0.7950 & 0.858 & -6.3 \\
    \midrule
    CWQ    & Hits@1   & 0.7241 & 0.817 & -9.3 \\
           & Macro-F1 & 0.7385 & 0.845 & -10.7 \\
           & EHR      & 0.7370 & 0.886 & -14.9 \\
    \bottomrule
  \end{tabular}
\end{table}

Three contributions to the CWQ gap were identified and quantified.

\paragraph{Superlatives are never learned.}
\num{1906} CWQ training questions (\SI{6.9}{\percent}) contain
\texttt{ORDER BY \ldots LIMIT}, but \texttt{gen\_parse\_data.py} does not parse
the sort clause, so the \texttt{order} field stays empty for all \num{27474}
records. Consequently \emph{no} training sample contains a \texttt{Rank} step,
although the prompt explicitly offers one. At test time \num{272} questions
(\SI{8.3}{\percent}) require a superlative; none of the reconstructed queries
contains \texttt{ORDER BY}. Without \texttt{LIMIT} the query returns all
candidates instead of the extremum, so recall stays high while precision
collapses -- these questions reach \(F_1 = 0.416\) against \(\approx 0.80\)
elsewhere. Repairing this is worth roughly \SI{3.2}{pp} of CWQ F1 but requires
regenerating the training data and retraining.

\paragraph{Converging graphs cannot be expressed.}
MemQ stores one parent per node (\texttt{all\_rel[node]['father']}), so a gold
graph in which two paths meet at the same variable is not representable. This
affects \num{53} of \num{3273} CWQ questions, which score \(F_1 = 0.145\)
against \(0.748\) for tree-shaped ones -- roughly \SI{1}{pp} overall.

\paragraph{The remainder is retrieval, not reconstruction.}
Of \num{2084} two-branch direction UNIONs, \num{805} have the first branch in
the gold set and \num{452} the second, so no simple rule would resolve the
direction. Crucially \num{827} (\SI{39.7}{\percent}) contain \emph{neither}:
the retrieved relation itself is wrong. The recall thresholds
\(\gamma_1, \gamma_2\) of Equation~4 are a plausible lever here -- the paper
defines the formula but states no values, so the 0.90/0.80 used here are
uncalibrated.

\section{Evaluation defects found and fixed}
\label{sec:defects}

Four defects were found while producing these numbers. All of them affected
\emph{every} model measured with this pipeline, not just v14.

\begin{table}[htbp]
  \centering
  \caption{Effect of the name-cache fix (v14 cp13750, WebQSP, \num{300}
  questions). Merging GrailQA entity names had overwritten \num{451} established
  WebQSP/CWQ names with GrailQA's shorter \texttt{friendly\_name} (e.g.\
  ``nixon'' for ``Richard Nixon''), so entity resolution failed for every
  affected question.}
  \label{tab:bugfix}
  \small
  \begin{tabular}{lS[table-format=2.4]S[table-format=2.4]}
    \toprule
    Metric & {Before fix} & {After fix} \\
    \midrule
    Reconstruction failures & 60 & 1 \\
    Structure accuracy      & 0.5033 & 0.6233 \\
    Hits@1                  & 0.6275 & 0.8020 \\
    Macro-F1                & 0.6273 & 0.7991 \\
    \bottomrule
  \end{tabular}
\end{table}

\paragraph{Whitespace.}
\SI{5.6}{\percent} of the DeepSeek memory descriptions begin with whitespace,
and the reconstruction unconditionally inserts a blank before every variable;
both produce double blanks, which the plan parser rejects. Initially
\emph{all} \num{300} reconstructions failed.

\paragraph{Structure accuracy blind to GrailQA.}
GrailQA's gold queries bind the Freebase namespace to \texttt{:} rather than
\texttt{ns:}, so the gold relation multiset came out empty and structure
accuracy was \num{0} by construction. Their type constraints additionally put a
class in object position, which matches the relation pattern and would be
compared against relations.

\paragraph{UNION branches invisible to EHR/GoldGED.}
\texttt{extract\_triples} split the WHERE block on \verb!" .\n"! and kept only
three-token lines. The reconstruction wraps a hop in
\verb!{ s p o }UNION{ o p s }! whenever the direction is unknown, so both
branches share a line, yield six tokens and were dropped -- every such edge
counted as missing. The symptom was diagnostic: EHR was worst on the
\emph{simplest} questions, which contain proportionally the most UNIONs, and
queries answering perfectly (\(F_1 = 1.0\)) scored \(\text{EHR} = 0\).

\begin{table}[htbp]
  \centering
  \caption{Edge hitting rate before and after correcting the triple
  extraction.}
  \label{tab:ehrfix}
  \begin{tabular}{lS[table-format=1.3]S[table-format=1.3]S[table-format=1.3]}
    \toprule
    Dataset & {Before} & {After} & {Paper} \\
    \midrule
    CWQ            & 0.497 & \bfseries 0.737 & 0.886 \\
    WebQSP         & 0.502 & \bfseries 0.795 & 0.858 \\
    WebQSP, 1 hop  & 0.437 & \bfseries 0.872 & 0.816 \\
    \bottomrule
  \end{tabular}
\end{table}

Answer-level metrics are unaffected by the last two fixes: they change how the
reconstructed graph is read back, not what is generated.

\section{Metric caveats}

\paragraph{Structure accuracy is definition-dependent.}
It is not a paper metric -- the paper defines only GoldGED (Eq.~5) and EHR
(Eq.~6). Because relation alternatives inflate the reconstructed multiset,
requiring exact equality penalises correct-but-broader queries.
Table~\ref{tab:sa} shows the spread; none of the variants is objectively
correct, so the figure should always be reported together with its definition.

\begin{table}[htbp]
  \centering
  \caption{Structure accuracy under three definitions (v9, full test sets).}
  \label{tab:sa}
  \small
  \begin{tabular}{lS[table-format=1.4]S[table-format=1.4]}
    \toprule
    Definition & {CWQ} & {WebQSP} \\
    \midrule
    exact multiset equality (current) & 0.2687 & 0.4440 \\
    one branch per UNION group        & 0.6270 & 0.7193 \\
    gold multiset contained           & 0.7194 & 0.8165 \\
    \bottomrule
  \end{tabular}
\end{table}

\paragraph{EHR as computed here is an upper bound.}
All UNION branches are treated as present edges, whereas the paper describes a
single reconstructed graph \(G_{re}\). Counting only the first branch per group
yields \num{0.658} for CWQ instead of \num{0.737}.

\paragraph{Hits@1 is not deterministic.}
Equation~7 defines it as \(\text{count}(\textit{Answer}[0] \in
\textit{Golden})\), and the implementation matches. However
\(\textit{Answer}[0]\) is the first row returned by the endpoint, and for
\texttt{SELECT DISTINCT} without \texttt{ORDER BY} that order is undefined in
SPARQL. Small run-to-run differences in Hits@1 are therefore expected, while
Macro-F1 is order-independent and reproduces exactly.

\end{document}
